\documentclass[10pt]{article}

% COMP 2300 Sets Core companion: student-facing reference and practice.
\usepackage[letterpaper,margin=0.56in,includehead,includefoot,
  headheight=15pt,headsep=7pt,footskip=20pt]{geometry}
\usepackage[T1]{fontenc}
\usepackage{amsmath,amssymb,mathtools}
\usepackage{microtype}
\usepackage[dvipsnames,svgnames]{xcolor}
\usepackage{multicol}
\usepackage{enumitem}
\usepackage{tabularx,booktabs,array}
\usepackage[most]{tcolorbox}
\usepackage{fancyhdr}
\usepackage{titlesec}
\usepackage[hidelinks]{hyperref}

\definecolor{CoreNavy}{HTML}{18364A}
\definecolor{CoreTeal}{HTML}{176B6B}
\definecolor{CoreGold}{HTML}{9A6500}
\definecolor{CoreRed}{HTML}{963D38}
\definecolor{CoreBlueTint}{HTML}{EDF4F8}
\definecolor{CoreTealTint}{HTML}{EDF8F6}
\definecolor{CoreGoldTint}{HTML}{FFF7E5}
\definecolor{CoreRedTint}{HTML}{FCEFED}

\setlength{\parindent}{0pt}
\setlength{\parskip}{2.5pt}
\setlength{\columnsep}{0.25in}
\setlength{\columnseprule}{0.25pt}
\setlength{\emergencystretch}{1.2em}
\setlength{\abovedisplayskip}{3pt}
\setlength{\belowdisplayskip}{3pt}
\setlist{leftmargin=*,itemsep=1pt,topsep=1.5pt,parsep=0pt,partopsep=0pt}
\raggedbottom

\titleformat{\section}{\Large\bfseries\color{CoreNavy}}{}{0pt}{}
  [\vspace{-3pt}\color{CoreNavy}\rule{\linewidth}{0.7pt}]
\titlespacing*{\section}{0pt}{3pt}{4pt}

\pagestyle{fancy}
\fancyhf{}
\fancyhead[L]{\textbf{COMP 2300}}
\fancyhead[C]{\textbf{Sets Core Companion}}
\fancyhead[R]{Definitions and Worked Examples}
\fancyfoot[L]{Core sections A--E}
\fancyfoot[C]{\thepage}
\fancyfoot[R]{Study reference---not used during a Core}
\renewcommand{\headrulewidth}{0.5pt}
\renewcommand{\footrulewidth}{0.3pt}
\makeatletter
\let\ps@plain\ps@fancy
\let\ps@empty\ps@fancy
\makeatother

\tcbset{enhanced,before skip=2.5pt,after skip=2.5pt,boxrule=0.55pt,arc=1.2mm,
  left=4pt,right=4pt,top=3pt,bottom=3pt,fonttitle=\bfseries}
\newtcolorbox{definitionbox}[1]{title={#1},colback=CoreBlueTint,
  colframe=CoreNavy,colbacktitle=CoreNavy,coltitle=white}
\newtcolorbox{methodbox}[1]{title={#1},colback=CoreTealTint,
  colframe=CoreTeal,colbacktitle=CoreTeal,coltitle=white}
\newtcolorbox{workedbox}[1]{title={#1},colback=white,
  colframe=CoreNavy,colbacktitle=CoreBlueTint,coltitle=CoreNavy}
\newtcolorbox{warningbox}[1]{title={#1},colback=CoreRedTint,
  colframe=CoreRed,colbacktitle=CoreRed,coltitle=white}
\newtcolorbox{checklistbox}[1]{title={#1},colback=CoreGoldTint,
  colframe=CoreGold,colbacktitle=CoreGold,coltitle=white}
\newcommand{\term}[1]{\textbf{\color{CoreNavy}#1}}
\newcommand{\U}{\mathfrak U}

\begin{document}
\small

\section*{A--B. Set Language, Constructions, and Regions}
\begin{multicols}{2}
\begin{definitionbox}{Elements, subsets, and equality}
$x\in A$ means $x$ is an element of $A$; $A\subseteq B$ means every element
of $A$ belongs to $B$. A proper subset $A\subset B$ also requires $A\ne B$.
Sets are equal exactly when $A\subseteq B$ and $B\subseteq A$. Order and
repetition do not matter inside a set.
\end{definitionbox}

\begin{definitionbox}{Roster and set-builder notation}
Roster notation lists elements; set-builder notation states a domain and a
property. For example,
\[
\{x\in\mathbb Z:-2\le x<3\}=\{-2,-1,0,1,2\}.
\]
Read the colon as ``such that.'' Always enforce the declared domain.
\end{definitionbox}

\begin{definitionbox}{Power sets and Cartesian products}
$\mathcal P(A)$ is the set of all subsets of $A$. If $|A|=n$, then
$|\mathcal P(A)|=2^n$. The empty set and $A$ itself are included.
\[
A\times B=\{(a,b):a\in A\text{ and }b\in B\}.
\]
Ordered pairs preserve order; generally $(a,b)\ne(b,a)$.
\end{definitionbox}

\begin{workedbox}{A worked construction}
Let $S=\{x\in\mathbb Z:|x|\le2\}$ and $T=\{-1,1\}$. Then
\[
S=\{-2,-1,0,1,2\},\quad T\subseteq S,\quad T\notin S,
\]
and
\[
\mathcal P(T)=\{\varnothing,\{-1\},\{1\},\{-1,1\}\}.
\]
\textbf{Step 1:} solve $|x|\le2$ in the declared integer domain.
\textbf{Step 2:} both elements of $T$ occur in $S$, so $T\subseteq S$.
\textbf{Step 3:} the elements of $S$ are integers, so the set $T$ is not itself
an element: $T\notin S$. \textbf{Step 4:} list subsets of $T$ by size
(zero, one, then two elements); the count $4=2^2$ checks the power set.
\end{workedbox}

\columnbreak

\begin{definitionbox}{Set operations}
Relative to universe $\U$,
\begin{align*}
A\cup B&=\{x:x\in A\text{ or }x\in B\},\\
A\cap B&=\{x:x\in A\text{ and }x\in B\},\\
A\setminus B&=\{x:x\in A\text{ and }x\notin B\},\\
\overline A&=\U\setminus A.
\end{align*}
Membership formulas and set formulas describe the same regions.
\end{definitionbox}

\begin{methodbox}{Translate logic into a Venn region}
Replace $x\in A$ by $A$, $\lor$ by union, $\land$ by intersection, and
$\neg$ by complement. First eliminate conditionals:
\[
P\to Q\equiv\neg P\lor Q,
\qquad
P\leftrightarrow Q\equiv(P\land Q)\lor(\neg P\land\neg Q).
\]
Then shade exactly the resulting region.
\end{methodbox}

\begin{workedbox}{B worked translations}
\begin{align*}
(x\in A\setminus B)\lor(x\in C)
&\longleftrightarrow (A\setminus B)\cup C,\\
(x\in A\cup B)\to(x\in C)
&\longleftrightarrow \overline{A\cup B}\cup C,\\
(x\in A)\leftrightarrow(x\in B\setminus C)
&\longleftrightarrow [A\cap(B\setminus C)]\\
&\quad\cup[\overline A\cap\overline{(B\setminus C)}].
\end{align*}
\textbf{Procedure:} (1) eliminate $\to$ or $\leftrightarrow$; (2) replace
$\neg,\land,\lor$ by complement, intersection, and union; (3) simplify.
For the biconditional, the first term is the region where both sides are true;
the second is where both are false. Shade both.
\end{workedbox}

\begin{warningbox}{Common errors}
$\{x\}\in\mathcal P(A)$ is equivalent to $x\in A$, but $\{x\}\in A$ is a
different claim. Include $\varnothing$ in every power set. Do not treat a
Cartesian product as unordered.
\end{warningbox}

\begin{checklistbox}{A--B readiness}
Can you move between roster and set-builder notation; distinguish $\in$ from
$\subseteq$; build $\mathcal P(A)$ and $A\times B$; and convert a compound
membership formula into a set expression and Venn region?
\end{checklistbox}
\end{multicols}

\clearpage
\section*{C. Relations as Subsets of Cartesian Products}
\begin{multicols}{2}
\begin{definitionbox}{Relations and plots}
A relation from $X$ to $Y$ is any subset $R\subseteq X\times Y$. A plotted
point at horizontal coordinate $x$ and vertical coordinate $y$ represents
$(x,y)\in R$. The \term{domain of the relation} is the set of first
coordinates that occur; its \term{range} is the set of second coordinates.
\end{definitionbox}

\begin{definitionbox}{When a relation is a function}
A relation $R\subseteq X\times Y$ is a function $X\to Y$ exactly when
\[
\forall x\in X\;\exists!y\in Y\;((x,y)\in R).
\]
Every input must have exactly one output. An input with no point or with two
points proves the relation is not a function.
\end{definitionbox}

\begin{definitionbox}{Injective and surjective}
For a function $f:X\to Y$:
\begin{align*}
\text{injective}&:\quad f(x_1)=f(x_2)\Rightarrow x_1=x_2,\\
\text{surjective}&:\quad \forall y\in Y\;\exists x\in X\;(f(x)=y),\\
\text{bijective}&:\quad \text{both injective and surjective.}
\end{align*}
On a finite plot, injective means no two inputs share an output; surjective
means every codomain value is hit.
\end{definitionbox}

\begin{workedbox}{C worked plot-to-pairs example}
Let $X=Y=\{1,2,3\}$ and suppose the plot has points
\[
R=\{(1,2),(2,3),(3,2)\}.
\]
\textbf{Function check:} columns $x=1,2,3$ each contain exactly one point, so
$R$ is a function $X\to Y$. \textbf{Injective check:} inputs $1$ and $3$ both
map to $2$, so it is not injective. \textbf{Surjective check:} codomain value
$1\in Y$ is never hit, so it is not surjective.
\end{workedbox}

\columnbreak

\begin{methodbox}{Read quantified plot statements}
\begin{align*}
\forall x\exists y\,R(x,y)&:\ \text{every column has a point},\\
\forall x\exists!y\,R(x,y)&:\ \text{every column has exactly one},\\
\forall y\exists x\,R(x,y)&:\ \text{every row has a point},\\
\forall y\exists!x\,R(x,y)&:\ \text{every row has exactly one},\\
\exists x\forall y\,R(x,y)&:\ \text{some column is full},\\
\exists y\forall x\,R(x,y)&:\ \text{some row is full}.
\end{align*}
The first variable is the horizontal coordinate in $(x,y)$.
\end{methodbox}

\begin{workedbox}{C quantified evaluation}
For $R=\{(1,2),(2,3),(3,2)\}$:
\begin{itemize}
\item $\forall x\exists!y\,R(x,y)$: scan every column; each has one point, so true.
\item $\forall y\exists x\,R(x,y)$: row $y=1$ is empty, a counterexample, so false.
\item $\exists y\forall x\,R(x,y)$: test each row; none contains all three
inputs, so false.
\end{itemize}
For a false universal, name the missing row or column; for a true existential,
name the row or column that is the witness.
\end{workedbox}

\begin{workedbox}{C nonfunction diagnosis}
If
$S=\{(1,1),(1,3),(2,2),(3,3)\}$, then input $1$ has two outputs. Thus $S$ is
not a function $X\to Y$, even though every input has at least one output.
Classification as injective or surjective is requested only after the relation
passes the function test.
\end{workedbox}

\begin{warningbox}{Common errors}
Do not confuse ``exactly one output for each input'' with ``exactly one input
for each output.'' A function may be many-to-one. Surjectivity is always
relative to the declared codomain, not merely the values visible in a formula.
\end{warningbox}

\begin{checklistbox}{C readiness}
Can you list a plotted relation as ordered pairs, evaluate each quantifier
pattern, identify a missing or duplicated output, and classify a valid function
as injective, surjective, both, or neither?
\end{checklistbox}
\end{multicols}

\clearpage
\thispagestyle{fancy}
\section*{D. Functions, Images, and Preimages}
\begin{multicols}{2}
\begin{definitionbox}{Domain, codomain, range}
In $f:X\to Y$, $X$ is the domain and $Y$ the codomain. The range is
$f(X)=\{f(x):x\in X\}$ and may be smaller than $Y$. The same formula can have
different injectivity or surjectivity behavior when its declared domain or
codomain changes.
\end{definitionbox}

\begin{definitionbox}{Images and preimages of sets}
For $U\subseteq X$ and $P\subseteq Y$,
\[
f(U)=\{f(x):x\in U\},
\qquad
f^{-1}(P)=\{x\in X:f(x)\in P\}.
\]
$f^{-1}(P)$ is preimage notation and does not require $f$ to have an inverse
function.
\end{definitionbox}

\begin{workedbox}{D finite-function example}
Let $A=\{1,2,3,4\}$, $B=\{a,b,c\}$, and
$f(1)=a,f(2)=b,f(3)=b,f(4)=c$. Then
\[
f(\{1,3,4\})=\{a,b,c\},
\quad
f^{-1}(\{a,b\})=\{1,2,3\}.
\]
\textbf{Image:} evaluate inputs $1,3,4$ to get $a,b,c$.
\textbf{Preimage:} collect all inputs whose outputs lie in $\{a,b\}$, namely
$1,2,3$. \textbf{Classification:} every element of $B$ is hit, so $f$ is
surjective; $f(2)=f(3)=b$ with $2\ne3$, so it is not injective.
\end{workedbox}

\begin{methodbox}{Proving or disproving injectivity}
To prove injective, start with arbitrary $x_1,x_2$ and assume
$f(x_1)=f(x_2)$; derive $x_1=x_2$. To disprove it, give distinct inputs with
the same output.
\end{methodbox}

\columnbreak

\begin{methodbox}{Proving or disproving surjectivity}
To prove surjective, let arbitrary $y$ lie in the codomain and solve
$f(x)=y$ for an allowed $x$ in the domain. To disprove it, identify one
codomain element with no preimage.
\end{methodbox}

\begin{workedbox}{D integer formula examples}
For $g:\mathbb Z\to\mathbb Z$, $g(n)=n+5$:
\[
g(\{-1,0,2\})=\{4,5,7\},\quad
g^{-1}(\{4,7\})=\{-1,2\}.
\]
\textbf{Injective:} assume $g(n_1)=g(n_2)$; then
$n_1+5=n_2+5$, hence $n_1=n_2$. \textbf{Surjective:} let arbitrary
$y\in\mathbb Z$; choosing $n=y-5\in\mathbb Z$ gives
$g(n)=(y-5)+5=y$. Thus $g$ is bijective.
\end{workedbox}

\begin{workedbox}{Domain and codomain matter}
For $h:\mathbb Z\to\mathbb Z$, $h(n)=n^2$, the inputs $-1$ and $1$ collide,
so $h$ is not injective. No integer squares to $-1$, so it is not surjective.
For $q:\mathbb N_0\to\{m^2:m\in\mathbb N_0\}$, $q(n)=n^2$, the domain removes
the $n,-n$ collisions. Every declared codomain value has the form $m^2=q(m)$
for an allowed input. Thus $q$ is injective and surjective. The formula did not
change; the declared domain and codomain did.
\end{workedbox}

\begin{warningbox}{Common errors}
An image removes duplicate outputs. A preimage collects inputs, not output
values. Never classify a function without reading both domain and codomain.
Do not claim surjectivity by saying ``the formula has outputs''; construct a
preimage for an arbitrary codomain element.
\end{warningbox}

\begin{checklistbox}{D readiness}
Can you compute images and preimages; distinguish range from codomain; prove or
disprove injectivity and surjectivity with the correct witnesses; and explain
how changing the domain or codomain changes the classification?
\end{checklistbox}
\end{multicols}

\clearpage
\section*{E. Prove or Refute Set and Function Claims}
\begin{multicols}{2}
\begin{methodbox}{Element-chasing proof template}
To prove $S=T$, prove both inclusions. For $S\subseteq T$:
\begin{enumerate}
\item Let $z\in S$ be arbitrary.
\item Unpack membership in $S$ using definitions.
\item Apply the relevant set/function definitions.
\item Conclude $z\in T$.
\end{enumerate}
Do not begin by assuming the desired equality.
\end{methodbox}

\begin{workedbox}{True claim: image of a union}
For $f:X\to Y$ and $U,V\subseteq X$,
$f(U\cup V)=f(U)\cup f(V)$.

\textbf{Forward inclusion:} if $y\in f(U\cup V)$, then $y=f(x)$ for some
$x\in U\cup V$. Thus
$x\in U$ or $x\in V$, so $y\in f(U)$ or $y\in f(V)$. This proves one
inclusion. \textbf{Reverse inclusion:} if $y\in f(U)\cup f(V)$, then $y=f(x)$ for some
$x\in U$ or some $x\in V$; in either case $x\in U\cup V$, so
$y\in f(U\cup V)$. Both inclusions prove equality.
\end{workedbox}

\begin{workedbox}{False claim: image of an intersection}
In general, $f(U\cap V)\ne f(U)\cap f(V)$. Let
$X=\{1,2\}$, $Y=\{a\}$, $f(1)=f(2)=a$, $U=\{1\}$, and $V=\{2\}$. Then
\[
f(U\cap V)=f(\varnothing)=\varnothing,
\quad
f(U)\cap f(V)=\{a\}.
\]
\textbf{Verification:} the hypotheses hold, but the left side is empty and the
right side is not, so the proposed universal equality is false. It fails because
two different inputs may have the same output.
\end{workedbox}

\columnbreak

\begin{workedbox}{True claim: preimage preserves intersection}
For $P,Q\subseteq Y$,
\[
f^{-1}(P\cap Q)=f^{-1}(P)\cap f^{-1}(Q).
\]
Indeed, for arbitrary $x\in X$,
\begin{align*}
x\in f^{-1}(P\cap Q)
&\Longleftrightarrow f(x)\in P\cap Q\\
&\Longleftrightarrow f(x)\in P\text{ and }f(x)\in Q\\
&\Longleftrightarrow x\in f^{-1}(P)\cap f^{-1}(Q).
\end{align*}
The chain proves both inclusions at once.
\end{workedbox}

\begin{workedbox}{True claim: product distributes over union}
For arbitrary $(x,y)$,
\begin{align*}
(x,y)\in(A\cup B)\times C
&\Longleftrightarrow (x\in A\text{ or }x\in B)\text{ and }y\in C\\
&\Longleftrightarrow [(x\in A\text{ and }y\in C)\text{ or}\\
&\hspace{3.4em}(x\in B\text{ and }y\in C)]\\
&\Longleftrightarrow (x,y)\in(A\times C)\cup(B\times C).
\end{align*}
\end{workedbox}

\begin{workedbox}{False cancellation claim}
$A\cup B=A\cup C$ need not imply $B=C$. Take
$A=\{1,2\}$, $B=\{1\}$, and $C=\{2\}$. Then both unions equal
$\{1,2\}$, so the hypothesis is true, but $B\ne C$, so the conclusion is
false. Union with $A$ hides differences already contained in $A$. A valid
counterexample must pass both of these checks.
\end{workedbox}

\begin{warningbox}{Common errors}
Examples do not prove universal claims. For a false equality, compute both
sides of a concrete finite example. In an image proof, different witnesses may
map to the same output; do not silently identify them.
\end{warningbox}

\begin{checklistbox}{Whole-Core readiness}
Can you construct finite sets, translate Venn regions, interpret relations,
compute and classify functions, and choose between an arbitrary-element proof
and a concrete counterexample? A live 20-minute form selects one complete
problem from each section A--E.
\end{checklistbox}
\end{multicols}

\end{document}
